\Chapter{6}

\Verse{1}
{
	\hRJE{}
}
{
	\eRJE{
		6

¹ And YHWH said to Moses, “Now you’ll see what I shall do to Pharaoh,
because with a strong hand he’ll let them go, and with a strong hand
he’ll drive them from his land.”\\
	}
}
{
}

\Verse{2}
{
	\hP{}
}
{
	\eP{
		AND I APPEARED

² And God spoke to Moses and said to him, “I am YHWH.\\
	}
}
{
}

\Verse{3}
{
	\hP{}
}
{
	\eP{
		And I appeared to Abraham, to Isaac, and to Jacob as El Shadday, and I
was not known to them by my name, YHWH.\\
	}
}
{
}

\Verse{4}
{
	\hP{}
}
{
	\eP{
		And I also established my covenant with them, to give them the land of
Canaan, the land of their residences, in which they resided.\\
	}
}
{
}

\Verse{5}
{
	\hP{}
}
{
	\eP{
		And also: I’ve heard the cry of the children of Israel as Egypt is
enslaving them, and I’ve remembered my covenant.\\
	}
}
{
}

\Verse{6}
{
	\hP{}
}
{
	\eP{
		Therefore, say to the children of Israel: ‘I am YHWH, and I shall bring
you out from under Egypt’s burdens, and I shall rescue you from their
toil, and I shall redeem you with an outstretched arm and with
tremendous judgments,\\
	}
}
{
}

\Verse{7}
{
	\hP{}
}
{
	\eP{
		and I shall take you to me as a people, and I shall become your God,
and you’ll know that I am YHWH, your God, who is bringing you out from
under Egypt’s burdens.\\
	}
}
{
}

\Verse{8}
{
	\hP{}
}
{
	\eP{
		And I shall bring you to the land that I raised my hand to give to
Abraham, to Isaac, and to Jacob, and I shall give it to you as a
possession. I am YHWH.’”\\
	}
}
{
}

\Verse{9}
{
	\hP{}
}
{
	\eP{
		And Moses spoke this to the children of Israel, and they did not listen
to Moses because of shortage of spirit and because of hard work.\\
	}
}
{
}

\Verse{10}
{
	\hP{}
}
{
	\eP{
		And YHWH spoke to Moses, saying,\\
	}
}
{
}

\Verse{11}
{
	\hP{}
}
{
	\eP{
		“Come, speak to Pharaoh, king of Egypt, that he should let the
children of Israel go from his land.”\\
	}
}
{
}

\Verse{12}
{
	\hP{}
}
{
	\eP{
		And Moses spoke in front of YHWH, saying, “Here, the children of
Israel didn’t listen to me, and how will Pharaoh listen to me?! And I’m
uncircumcised of lips!”

6:12,30. uncircumcised of lips. Does this shed light on the meaning of
“heavy of tongue and heavy of mouth”? It seems rather to be even harder
to interpret. Uncircumcision is used elsewhere in the Torah as a
metaphor as well. Lev 26:41 and Deut 10:16 refer to the people’s
uncircumcised heart. (The prophets Jeremiah and Ezekiel use this image
also.) There it is a strong image, criticizing the people who understand
that their covenant requires circumcision of the skin that covers the
penis, but who do not appreciate that they must likewise remove the
impediment to their hearts. But the meaning of the image here is less
clear. We still do not know what the impediment is. All we can say is
that Moses is again pictured as unconfident, feeling inadequate to the
task. The repeated emphasis on this point will make the eventual success
of the exodus that much more dramatic.\\
	}
}
{
}

\Verse{13}
{
	\hP{}
}
{
	\eP{
		And YHWH spoke to Moses and to Aaron and commanded them regarding the
children of Israel and regarding Pharaoh, king of Egypt, to bring out
the children of Israel from the land of Egypt.\\
	}
}
{
}

\Verse{14}
{
	\hP{}
}
{
	\eP{
		These are the heads of their fathers’ houses:

The sons of Reuben, Israel’s firstborn: Hanoch and Pallu, Hezron, and
Carmi. These are the families of Reuben.\\
	}
}
{
}

\Verse{15}
{
	\hP{}
}
{
	\eP{
		And the sons of Simeon: Jemuel and Jamin and Ohad and Jachin and Zohar
and Saul, son of a Canaanite woman. These are the families of Simeon.\\
	}
}
{
}

\Verse{16}
{
	\hP{}
}
{
	\eP{
		And these are the names of the sons of Levi by their records: Gershon
and Kohath and Merari. And the years of Levi’s life were a hundred
thirty-seven years.\\
	}
}
{
}

\Verse{17}
{
	\hP{}
}
{
	\eP{
		The sons of Gershon: Libni and Shimei, by their families.\\
	}
}
{
}

\Verse{18}
{
	\hP{}
}
{
	\eP{
		And the sons of Kohath: Amram and Izhar and Hebron and Uzziel. And the
years of Kohath’s life were a hundred thirty-three years.\\
	}
}
{
}

\Verse{19}
{
	\hP{}
}
{
	\eP{
		And the sons of Merari: Mahli and Mushi. These are the families of
Levi by their records.\\
	}
}
{
}

\Verse{20}
{
	\hP{}
}
{
	\eP{
		And Amram took Jochebed, his aunt, as a wife; and she gave birth to
Aaron and Moses for him. And the years of Amram’s life were a hundred
thirty-seven years.\\
	}
}
{
}

\Verse{21}
{
	\hP{}
}
{
	\eP{
		And the sons of Izhar: Korah and Nepheg and Zichri.\\
	}
}
{
}

\Verse{22}
{
	\hP{}
}
{
	\eP{
		And the sons of Uzziel: Mishael and Elzaphan and Sithri.\\
	}
}
{
}

\Verse{23}
{
	\hP{}
}
{
	\eP{
		And Aaron took Elisheba, daughter of Amminadab, sister of Nahshon as a
wife; and she gave birth to Nadab and Abihu, Eleazar, and Ithamar for
him.\\
	}
}
{
}

\Verse{24}
{
	\hP{}
}
{
	\eP{
		And the sons of Korah: Assir and Elkanah and Abiasaph. These are the
families of the Korahites.\\
	}
}
{
}

\Verse{25}
{
	\hP{}
}
{
	\eP{
		And Eleazar son of Aaron had taken one of the daughters of Putiel as a
wife, and she gave birth for him to Phinehas. These are the heads of the
fathers of the Levites according to their families.\\
	}
}
{
}

\Verse{26}
{
	\hP{}
}
{
	\eP{
		That is Aaron and Moses, to whom YHWH said, “Bring out the children of
Israel from the land of Egypt by their masses.”

6:26–27. Aaron and Moses … Moses and Aaron. Both word orders are given:
Aaron first, then Moses first, suggesting that each is important in a
different way. The recognition of Aaron here is especially significant
because it is natural to think of Moses first, but this notice follows a
genealogy that has listed Aaron as the firstborn and identified his wife
and children. (The Septuagint and Qumran texts have the order in v. 27
reversed, which unfortunately loses this point.)\\
	}
}
{
}

\Verse{27}
{
	\hP{}
}
{
	\eP{
		They were the ones speaking to Pharaoh, king of Egypt, to bring out
the children of Israel from Egypt. That is Moses and Aaron.\\
	}
}
{
}

\Verse{28}
{
	\hP{}
}
{
	\eP{
		And it was in the day that YHWH spoke to Moses in the land of Egypt,\\
	}
}
{
}

\Verse{29}
{
	\hP{}
}
{
	\eP{
		and YHWH spoke to Moses, saying, “I am YHWH. Speak to Pharaoh, king of
Egypt, everything that I speak to you.”\\
	}
}
{
}

\Verse{30}
{
	\hP{}
}
{
	\eP{
		And Moses said in front of YHWH, “Here, I’m uncircumcised of lips, and
how will Pharaoh listen to me?!”\\
	}
}
{
}

